\scanpage{16}

ne peut suffire à réconstruire $Q$ à isom. unique près -- mais seulement $Q$ ``modulo antipodisme''. Considérant le carré épinglé $Q_0$ et le groupe de ses automorphismes
\[
\deleted{\mathcal D_4}\ \mathbb D_4=\operatorname{Aut}(Q_0)
\]
avec la suite exacte canonique
\diagram{\begin{tikzcd}[ampersand replacement=\&,column sep=small]
1\arrow[r]\&\mathbb F_2\arrow[r]\&\mathbb D_4\arrow[r]\&\mathbb F_2^2\arrow[r,dash]\&1
\end{tikzcd}}
et la donnée de notre $\mathcal D$ revient à celle d'un torseur sous \inserted{le quotient} $\mathbb F_2^2$ \inserted{de $\mathbb D_4$} i.e. de deux ens. à deux éléments (ce sont en fait $D$ et $C$, dans cet ordre \ldots), tandis que celle d'un carré revient à la donnée d'un torseur sur $\mathbb D_4$ lui-$=$.

On a une opération \inserted{(par autom. int.)} de $\mathcal D^+$ sur $D\inserted{\hookrightarrow\mathcal D^+}$, la donnée de $Q$ (correspondant au groupe d'autom $\mathcal D$) \deleted{\unsure{ne}} revient à celle d'un relèvement de \deleted{\unsure{cette}} \deleted{\inserted{\unsure{cet}}} l'\deleted{\unsure{é}}opération \deleted{\unsure{induite de}} en \deleted{\unsure{$\mathcal D^+$|des|les}} une opération \deleted{\unsure{sur}} transitive sur un ens. $S$ \inserted{(donc un esp. hom sous $\mathcal D$)}, avec $S/\underline a\simeq D$. Description duale via $C\hookrightarrow\mathcal D^+\ldots$

\uline{Remarque.} Le groupe $\operatorname{Aut}(\mathbb D_4)$ est isom. à $\mathbb D_4$ (mais \uline{pas} via autom. intérieurs, bien sûr) ce qui implique qu'il y a \inserted{quand même} une équivalence de groupoïdes $(\text{Carrés})\xrightarrow{\approx}(\text{groupes isom. à }\mathbb D_4)$ (mais qui n'est \unsure{pas|$\vDash$} donnée par \deleted{\unsure{$\mathbb D_4$)}} $Q\mapsto\operatorname{Aut}(Q)$ !), qui mériterait d'être comprise \ldots

% archive sheet number: 15

